\chapter{Objetivos y metodología de trabajo}
\label{cap:objetivos}

% GUION. CAPITULO NUEVO (ago-2026), exigido por Instrucciones para la redaccion del TFE
% seccion 2.5. Los tres elementos son OBLIGATORIOS: (1) objetivo general, (2) objetivos
% especificos, (3) metodologia de trabajo.
%
% ⚠️ Los objetivos DEBEN ser SMART. La norma lo llama "imprescindible":
%      S especifico · M medible · A alcanzable · R relevante · T con plazo
%    Un objetivo sin criterio de exito medible es el fallo que la rubrica penaliza en
%    "Alcance" (10 %) y en "Relacion entre objetivos, planteamiento, desarrollo y
%    conclusiones" (10 %).
%
% ⚠️ Los objetivos que habia en el capitulo 1 eran del alcance antiguo (hablaban de
%    mitigacion de lectura y de validacion en hardware real). El alcance vigente,
%    acordado con el director, es el modulo GEM planteado como COMPARATIVA DE TRES
%    MODELOS — tipo de trabajo 3 de la titulacion.

\section{Objetivo general}
\label{sec:objetivo_general}
% Redactar en UNA frase que contenga el criterio de exito medible y el plazo.
% BORRADOR PROPUESTO — pendiente de confirmar con el director:
%
%   «Determinar, en el transcurso del presente curso academico, si la estructura de un
%    circuito cuantico combinada con la telemetria real del procesador permite predecir
%    la degradacion por ruido antes de ejecutarlo, evaluando el resultado sobre un
%    conjunto de test fuera de distribucion en tres ejes y considerando el objetivo
%    alcanzado si al menos uno de los tres modelos comparados reduce el error absoluto
%    medio frente a no mitigar.»
%
% Desglose SMART para justificarlo en el texto:
%   S — predecir la degradacion pre-ejecucion a partir de circuito + calibracion.
%   M — mejora del MAE frente al baseline de NO MITIGAR, medida en test OOD.
%   A — el conjunto de datos ya esta generado y los tres modelos son estandar.
%   R — sin mitigacion, los resultados de un procesador NISQ no son utilizables.
%   T — un cuatrimestre, el que dura la asignatura.

\section{Objetivos específicos}
\label{sec:objetivos_especificos}
% Verbos tomados de la lista aprobada por la norma: Construir/Definir/Comparar/
% Cuantificar/Determinar/Medir/Verificar. Cada uno con su criterio de exito.
\begin{enumerate}
    \item Construir un conjunto de datos de circuitos cuánticos con calibración real y
          tres ejes de generalización independientes —tipo de circuito, tamaño y tiempo—
          verificando que ningún día de calibración se comparte entre particiones.
    \item Definir una representación del circuito que no dependa del número de qubits,
          comprobando que el mismo modelo acepta circuitos de tamaño no visto en
          entrenamiento.
    \item Determinar qué magnitud del error es predecible antes de ejecutar, midiendo el
          $R^2$ en validación sobre cada candidata a variable objetivo.
    \item Comparar Ridge, Random Forest y un \textit{Graph Transformer} bajo un protocolo
          idéntico, reportando las métricas desglosadas por los tres ejes.
    \item Cuantificar qué aporta la estructura del circuito frente a un resumen agregado,
          mediante la diferencia de error entre el modelo de grafo y los tabulares.
\end{enumerate}

% ⚠️ Cada uno de estos objetivos debe cerrarse con una conclusion en el capitulo 12.
%    La norma lo pide explicitamente: «cada objetivo del trabajo se enlazara con una
%    conclusion».

\section{Metodología de trabajo}
\label{sec:metodologia}
% La norma pide: que pasos se dan, el PORQUE de cada paso, que instrumentos se usan y
% como se analizan los resultados.

\subsection{Tipo de trabajo y encaje en la titulación}
\label{sec:tipo_trabajo}
% Declararlo explicitamente: el tribunal lee los capitulos contra la estructura esperada
% del tipo elegido.
Comparativa de soluciones (tipo 3). Líneas de trabajo: aprendizaje automático —Ridge y
Random Forest— y aprendizaje profundo —el \textit{Graph Transformer}—.

\subsection{Fases del trabajo}
\label{sec:fases}
% Una fase por bloque de capitulos, para que se vea la trazabilidad con el resto de la
% memoria. Enunciar el porque de cada fase, no solo el que.
Revisión de la literatura y localización del hueco · diseño y generación del conjunto de
datos · análisis exploratorio y validación de las etiquetas · determinación de la variable
objetivo · entrenamiento de los tres modelos bajo protocolo común · análisis comparativo.

\subsection{Herramientas empleadas}
\label{sec:herramientas}
% La norma pide declarar QUE INSTRUMENTOS se utilizan. Aqui van TODOS, en pie de
% igualdad: el stack tecnico y las herramientas de IA.
Simulación y modelo de ruido, entrenamiento de los modelos, seguimiento de experimentos e
infraestructura de cómputo. Citar cada herramienta en formato APA.

\subsubsection{Declaración de uso de herramientas de inteligencia artificial}
\label{sec:declaracion_ia}
% ⚠️ OBLIGATORIO. Guia del TFE seccion 4.4: el uso «debe quedar explicitamente reflejado
%    en el documento incluyendo las referencias especificas a las herramientas de IA
%    utilizadas segun las normas de citacion oportunas», y «el director de TFE debe
%    validar su uso adecuado y etico».
%    Sancion en caso de uso no declarado o mal citado: NO AUTORIZAR el trabajo y
%    declararlo no apto para la defensa, con perdida de la convocatoria.
%
% ⚠️ La declaracion debe ser COMPLETA. Alcance acordado con el usuario (ago-2026):
%      - Claude Code — desarrollo del codigo de simulacion con Qiskit y de todo lo
%        ajeno al aprendizaje automatico;
%      - Claude Code — generacion de las figuras del trabajo;
%      - Claude Code — asistencia en la redaccion, estructuracion y revision de esta
%        memoria;
%      - Gemini — busqueda bibliografica profunda para el estado del arte.
%
% ⚠️ Pendiente: confirmar la validacion con el director y añadir las entradas APA de
%    cada herramienta a bibliografia.bib.
Declarar aquí, con detalle y sin ambigüedad, qué herramienta se usó para qué.
